\clearpage
\section{A rational truncation that fails the output contact test}\label{sec:truncation}
\subsection{A genuine algebraic candidate can be tested exactly}
Fix integers $r\ge1$ and $K\ge r$. Replace each polylogarithm by its finite series through index $K$, obtaining a rational polynomial $\mathcal H_{r,K}(t)$. Its Cauchy--Binet expansion has the same positive subset coefficients, restricted to $1\le k_i\le K$.

Let
\[
 L_K=\operatorname{lcm}(1,\ldots,K),\quad s_r=r(r+1)/2,
 \quad d=rK-r(r-1)/2.
\]
Then
\[
 A_{r,K}(t):=L_K^{r^2}\mathcal H_{r,K}(t)\in\Z[t],
 \qquad \ord_0A_{r,K}=s_r,\qquad \deg A_{r,K}=d.
\]
To see integrality, scale row $i$ by $L_K^{i+1}$ and column $j$ by $L_K^j$. Each entry becomes
\[
 L_K^{i+j+1}\sum_{k=1}^K\frac{t^k}{k^{i+j+1}}\in\Z[t].
\]
The total determinant scaling exponent is $r^2$. The unique smallest subset is $(1,\ldots,r)$ and the unique largest is $(K-r+1,\ldots,K)$, proving the order and degree.

For a reduced rational argument $t=a/c$, an obvious integral scalar is
\[
 W_{r,K}(a,c)=c^dA_{r,K}(a/c).
\]
It is nonzero when $0<a<c$, since the truncated positive expansion is nonzero. Thus it is a genuine nonzero arithmetic candidate. The question is whether it retains the required output-prime contact after clearing denominators.

\begin{theorem}[Loss of output contact]\label{thm:truncation-failure}
If $p>K$, $p\mid c$ and $\gcd(a,c)=1$, then
\[
 p\nmid W_{r,K}(a,c).
\]
\end{theorem}
\begin{proof}
Expand $W$ as a homogeneous polynomial in $a,c$. Modulo $p$, every term below the top degree in $a$ retains a positive power of $c$ and vanishes. The unique top subset is $(K-r+1,\ldots,K)$. Its indices, nonzero differences and $L_K$ are all $p$-units, so its scaled coefficient is a unit. The remaining factor $a^d$ is a unit because $p\nmid a$. The surviving term is therefore nonzero modulo $p$.
\end{proof}
If an output prime above the cutoff has repeated depth, this witness has no positive valuation there at all. The failed candidate does not pass the local output gate. It is not a theorem excluding every possible determinant realization.

\subsection{Rank two exposes the loss in one line}
The first two truncated rank-two expressions are
\[
 \mathcal H_{2,2}(t)=t^3/8,
\]
\[
 \mathcal H_{2,3}(t)=\frac{t^3}{8}+\frac{4t^4}{27}+\frac{t^5}{216}
 =\frac{t^3(27+32t+t^2)}{216}.
\]
Clearing the displayed denominator and the projective denominator in the second gives
\[
 W(a,c)=a^3(27c^2+32ac+a^2).
\]
For every $p>3$ dividing $c$, this is congruent to $a^5$ modulo $p$. Hence it is a unit at that output prime.

The original local series had contact near zero. The algebraic scalar evaluates at $a/c$ and clears powers of $c$ dictated by its top degree. At primes dividing $c$, it is the top term that survives. Long truncation does not automatically solve this mismatch. A valid construction must prove the contact at the actual local chart and retain it through the exact normalization map.
